\section{Discussion}\label{sec:discussion}

Our findings reveal both the strengths and boundaries of IRT person-fit for contamination detection: $\ell_z^*$ reliably detects SFT-stage contamination across model families, but three boundary conditions constrain its practical utility.

\paragraph{Characterization over tool.}
The detection-diagnosis tradeoff (\S\ref{sec:degradation}) shows that simpler statistics detect contamination more powerfully than $\ell_z^*$: accuracy $z$-scores yield $|d|$ values several-fold to over an order of magnitude larger, and HT generally exceeds $\ell_z^*$.
The value of the IRT framework lies not in detection power but in mechanistic understanding: the $\theta$-shift analysis reveals \emph{why} detection sensitivity varies across models and contamination levels, the seen/unseen decomposition identifies \emph{what fraction} of items are memorized, and the 4PL parameterization captures item-level ceiling effects that 3PL misattributes as aberrance.
These diagnostics are invisible to accuracy $z$-scores or HT; our contribution lies in quantifying these theoretically predictable effects in the LLM setting, where magnitudes and boundary conditions are not derivable from theory alone.

\paragraph{Practical implications.}
Two results have immediate consequences.
Free-$\theta$ estimation produces seed-dependent dose-response curves (CV up to 60\%); $\theta$-anchoring stabilizes detection for models whose clean $\ell_z^*$ departs substantially from zero.
$\ell_z^*$ is most informative for large benchmarks ($n \approx 8{,}000$ items for 80\% power at 15\% contamination); on small benchmarks like ARC-C, detection relies on concentrated contamination effects.

\paragraph{Scope and calibration.}
Universal baseline misfit (57/60 cells flagged) confirms systematic IRT violations; the $\Delta\ell_z^*$ design turns this into a feature through within-model differencing.
PSN-IRT provides pre-calibrated 4PL parameters \citep{psnirt2026}, but classical calibration fails on contaminated data ($r_a{=}0.127$; Appendix~\ref{app:mml_em}); the robustness analysis (Appendix~\ref{app:param_perturb}) provides partial reassurance (${\pm}10\%$ noise tolerance).
\paragraph{Information requirements and method positioning.}
$\ell_z^*$ sits below MIA \citep[logits;][]{minkprob2024} and ConStat \citep[multi-benchmark scores;][]{constat2024} in an information-requirements hierarchy, needing only single-benchmark binary patterns and an external item bank---suited to retrospective audits where logits are unavailable. Our GRPO pilot confirms this: Min-K\% AUC drops to chance ($0.493$) post-RL while $\ell_z^*$ signal persists (${\sim}3.3$; Figure~\ref{fig:grpo}), consistent with \citet{wang2025fragility}.
